\chapter{Nonlinear Controllability in Depth}
\label{ch:nonlinear-controllability}

\section{Scope of This Chapter}
\label{sec:controllability-scope}

\cref{ch:wmr-mechanics} introduced the controllability test as an
operational recipe: build the accessibility distribution with the Lie
bracket \cref{eq:lie-bracket}, check the rank condition \cref{eq:arc},
and read the degree of nonholonomy \( \kappa \)
(\cref{subsec:controllability}). The kinematic models, their
controllability, and the chained-form transformations, with the
related computations, are in that chapter.

This chapter does \emph{not} repeat those definitions and results. It
answers the question that the recipe leaves open: \emph{why does it
work?} It collects

\begin{itemize}
    \item the formal statement of the controllability criterion and an
          outline of its proof (\cref{sec:controllability-problem});
    \item the properties of the Lie bracket and the derivation of its
          geometric meaning (\cref{sec:lie-brackets,%
          sec:lie-bracket-interpretation});
    \item the link between involutivity, integrability and holonomy
          (\cref{sec:distributions});
    \item exercises with solutions, including the bracket
          computations for the vehicles of \cref{ch:wmr-mechanics}
          (\cref{sec:controllability-exercises}).
\end{itemize}

Throughout, the state is denoted by \( x \in \mathbb{R}^n \) when
talking about a generic system, and by \( q \) for the vehicles. The
system is the driftless one of \cref{eq:kinematic-model},
\begin{equation}
\label{eq:driftless-system}
    \dot{x} = \sum_{i=1}^{m} g_i(x)\, u_i .
\end{equation}

\section{Controllability of Driftless Systems}
\label{sec:controllability-problem}

\begin{definition}[Controllability of a Nonlinear System]
\label{def:controllability-nonlinear}
The system \cref{eq:driftless-system} is \textbf{controllable} if, for
any pair of states \( (x_i, x_f) \), with \( x(0) = x_i \), there
exists a profile \( \tilde{u}(t) \) for the inputs on
\( t \in [0, T] \) such that the system reaches \( x(T) = x_f \).
\end{definition}

For linear systems the test is a rank condition on the controllability
matrix. For nonlinear systems the test is again a rank condition, but
on a matrix built from Lie brackets, and the rank condition is, in
general, only a statement about \emph{accessibility}, i.e.\ about the
reachable set having non-empty interior. Driftless systems are special:
the two notions coincide. This is the content of the following theorem
(stated here, and used in \cref{ch:wmr-mechanics}, as the
accessibility rank condition).

\begin{theorem}[Accessibility Criterion (Chow--Rashevskii)]
\label{thm:accessibility}
Consider the driftless system \cref{eq:driftless-system} on a
connected domain, and let \( \Delta_{\mathcal{A}} \) be its
accessibility distribution (\cref{subsec:controllability}).
\begin{enumerate}
    \item If \( \dim \Delta_{\mathcal{A}}(x) = n \) at every \( x \),
          the system is controllable.
    \item Conversely, if the system is analytic and
          \( \dim \Delta_{\mathcal{A}} = \nu < n \), it is not
          controllable: from any initial state the reachable set is
          contained in a \( \nu \)-dimensional sub-manifold.
\end{enumerate}
\end{theorem}

\begin{nb}[Drift]
If the system has a drift term, \( \dot{x} = f(x) + \sum_i g_i(x)u_i \),
the rank condition on \( \Delta_{\mathcal{A}} \) only guarantees
accessibility: one cannot reverse the motion along \( f \), so the
reachable set need not be a neighbourhood. The wheeled-robot kinematic
models in this course are driftless, so the criterion above is a
necessary and sufficient test.
\end{nb}

\paragraph{Outline of the proof.}
The proof of part 1 has three steps.
\begin{enumerate}
    \item \emph{Brackets are reachable directions.} By
          \cref{sec:lie-bracket-interpretation}, a cyclic input
          sequence of duration \( 4\varepsilon \) produces a
          displacement \( \varepsilon^2 [g_i, g_j] + O(\varepsilon^3) \).
          Iterating the construction gives the same for nested
          brackets, with displacement of order \( \varepsilon^{k} \) for
          a bracket of length \( k \).
    \item \emph{Full rank gives an open set.} If the vector fields and
          their brackets span \( \mathbb{R}^n \) at \( x_0 \), suitable
          combinations of such sequences move the state in \( n \)
          independent directions, so the reachable set from \( x_0 \)
          contains an open neighbourhood of \( x_0 \).
    \item \emph{Globalization.} The system is driftless, so every
          motion can be reversed with the opposite inputs: ``being
          reachable from'' is an equivalence relation. Its classes are
          open by step 2, hence they also are closed (the complement is
          a union of open classes). On a connected domain there is
          only one class: the whole domain.
\end{enumerate}
Part 2 follows from the Frobenius theorem applied to the (involutive)
distribution \( \Delta_{\mathcal{A}} \), see \cref{sec:distributions}.

\paragraph{Holonomic vs.\ nonholonomic, revisited.}
Combining \cref{thm:accessibility} with the classification of
\cref{subsec:controllability}: the kinematic model of a robot is
controllable if and only if its kinematic constraints are
\emph{completely} nonholonomic (\( \nu = n \)). Partially integrable
constraints are still nonholonomic, but the system is not controllable.

\section{Lie Brackets in Depth}
\label{sec:lie-brackets}

The definition of the Lie bracket is \cref{eq:lie-bracket}. Using
the Jacobians \( J_i(x) = \partial g_i / \partial x \), it reads
\( [g_1, g_2] = J_2 g_1 - J_1 g_2 \).

\paragraph{Algebraic properties.}
For vector fields \( f, g, h \) and constants \( a, b \):
\begin{align}
    [f, g] &= -[g, f] && \text{(antisymmetry; in particular } [f,f]=0\text{)}, \\
    [af + bg, h] &= a[f, h] + b[g, h] && \text{(bilinearity)}, \\
    [f,[g,h]] + [g,[h,f]] + [h,[f,g]] &= 0 && \text{(Jacobi identity)}.
\end{align}
Antisymmetry is immediate from the definition and is the reason why, in
the accessibility distribution, each pair \( g_i, g_j \) is counted only
once. Bilinearity follows from the linearity of the Jacobian. The Jacobi
identity is verified by expanding the six terms: the second-derivative
terms cancel in pairs by the symmetry of the Hessian. Its practical
consequence is that, to span all brackets of a given length, it is
enough to consider \emph{left-nested} brackets
\( [g_{i_1}, [g_{i_2}, [\dots, g_{i_k}]]] \).

\paragraph{Example in \( \mathbb{R}^3 \).}
Let
\[
g_1 = \begin{pmatrix} \cos x_3 \\ \sin x_3 \\ 0 \end{pmatrix},
\qquad
g_2 = \begin{pmatrix} 0 \\ 0 \\ x_1 \end{pmatrix}.
\]
The Jacobians are
\[
J_1 =
\begin{pmatrix}
    0 & 0 & -\sin x_3 \\
    0 & 0 & \cos x_3 \\
    0 & 0 & 0
\end{pmatrix},
\qquad
J_2 =
\begin{pmatrix}
    0 & 0 & 0 \\
    0 & 0 & 0 \\
    1 & 0 & 0
\end{pmatrix},
\]
hence
\[
[g_1, g_2](x) = J_2 g_1 - J_1 g_2 =
\begin{pmatrix} 0 \\ 0 \\ \cos x_3 \end{pmatrix}
- \begin{pmatrix} -x_1\sin x_3 \\ x_1\cos x_3 \\ 0 \end{pmatrix}
=
\begin{pmatrix}
    x_1 \sin x_3 \\
    -x_1 \cos x_3 \\
    \cos x_3
\end{pmatrix}.
\]
The determinant of \( (g_1, g_2, [g_1,g_2]) \) is \( x_1^2 \), so
\( [g_1, g_2] \) is linearly independent from \( g_1, g_2 \) only for
\( x_1 \neq 0 \): on the plane \( x_1 = 0 \) the rank condition fails
at the first bracket, and one would have to compute higher-order
brackets.

\section{Interpretation of the Lie Bracket}
\label{sec:lie-bracket-interpretation}

Consider the 2-input system \( \dot{x} = g_1(x) u_1 + g_2(x) u_2 \)
and apply the following input profile on \( t \in [0, 4\varepsilon] \):
\begin{align}
    (u_1, u_2) &= (1, 0) & t &\in [0, \varepsilon], \\
    (u_1, u_2) &= (0, 1) & t &\in [\varepsilon, 2\varepsilon], \\
    (u_1, u_2) &= (-1, 0) & t &\in [2\varepsilon, 3\varepsilon], \\
    (u_1, u_2) &= (0, -1) & t &\in [3\varepsilon, 4\varepsilon].
\end{align}

\begin{figure}[ht]
    \centering
    \includegraphics[width=0.6\linewidth]{images/lie_inputs.eps}
    \caption{Input profile for the interpretation of the Lie bracket.}
    \label{fig:lie-inputs}
\end{figure}

\begin{proposition}[Wriggle]
\label{prop:wriggle}
After this input sequence,
\begin{equation}
\label{eq:wriggle}
    x(4\varepsilon) = x(0) + \varepsilon^2 [g_1, g_2](x(0)) + O(\varepsilon^3).
\end{equation}
\end{proposition}

\begin{proof}
The flow of a vector field \( g \) for a time \( t \) is, to second
order,
\[
\phi^{g}_t(x) = x + t\, g(x) + \tfrac{t^2}{2} J_g(x)\, g(x) + O(t^3),
\]
and the flow for \( -g \) is obtained with \( t \to -t \). Write
\( g_i, J_i \) for the fields and Jacobians evaluated at \( x(0) \),
and follow the four phases.
\begin{align*}
x_1 &= x + \varepsilon g_1 + \tfrac{\varepsilon^2}{2} J_1 g_1, \\
x_2 &= x_1 + \varepsilon g_2(x_1) + \tfrac{\varepsilon^2}{2} J_2 g_2
     = x + \varepsilon (g_1 + g_2)
       + \varepsilon^2 \big( \tfrac12 J_1 g_1 + J_2 g_1 + \tfrac12 J_2 g_2 \big),
\end{align*}
using \( g_2(x_1) = g_2 + \varepsilon J_2 g_1 + O(\varepsilon^2) \).
In the third phase,
\( g_1(x_2) = g_1 + \varepsilon J_1 (g_1 + g_2) + O(\varepsilon^2) \), so
\begin{align*}
x_3 &= x_2 - \varepsilon g_1(x_2) + \tfrac{\varepsilon^2}{2} J_1 g_1 \\
    &= x + \varepsilon g_2
       + \varepsilon^2 \big( J_2 g_1 - J_1 g_2 + \tfrac12 J_2 g_2 \big),
\end{align*}
where the \( J_1 g_1 \) terms cancel (\( \tfrac12 - 1 + \tfrac12 = 0 \)).
Finally, with \( g_2(x_3) = g_2 + \varepsilon J_2 g_2 + O(\varepsilon^2) \),
\[
x_4 = x_3 - \varepsilon g_2(x_3) + \tfrac{\varepsilon^2}{2} J_2 g_2
    = x + \varepsilon^2 (J_2 g_1 - J_1 g_2) + O(\varepsilon^3)
    = x + \varepsilon^2 [g_1, g_2] + O(\varepsilon^3). \qedhere
\]
\end{proof}

The displacement is in the direction of the Lie bracket, with an error
\( O(\varepsilon^3) \) that becomes negligible as
\( \varepsilon \to 0 \). A suitable cyclic sequence of inputs (a
\emph{wriggle}) therefore produces a net motion along a direction that
is not directly spanned by the inputs. Note the scaling:
the displacement is of order \( \varepsilon^2 \) for a maneuver of
duration \( 4\varepsilon \), so motion along a bracket is
``slower'' than motion along the input fields, and brackets of length
\( k \) are slower still (order \( \varepsilon^k \)).

\begin{figure}[ht]
    \centering
    \includegraphics[width=\linewidth]{images/lie_bracket_direction.eps}
    \caption{The Lie bracket gives the direction of the net displacement
    produced by the cyclic input sequence.}
    \label{fig:lie-bracket-direction}
\end{figure}

\begin{nb}[A Numerical Check of the Wriggle]
For the unicycle at \( \theta = 0 \), apply the sequence
\( (v, \omega) = (1, 0), (0, 1), (-1, 0), (0, -1) \), each for a time
\( \varepsilon \). The vehicle moves forward by \( \varepsilon \),
turns left by \( \varepsilon \), moves backward by \( \varepsilon \)
along the new heading, and turns back. The final position is
\( (\varepsilon - \varepsilon\cos\varepsilon,\
-\varepsilon\sin\varepsilon) \approx (0, -\varepsilon^2) \) with
\( \theta = 0 \), i.e.\ \( \varepsilon^2 [g_1, g_2] \) with
\( [g_1, g_2] = (\sin\theta, -\cos\theta, 0)^\top \) the slide vector
field of \cref{ch:wmr-mechanics}, in agreement with \cref{eq:wriggle}:
a lateral displacement obtained with no lateral velocity.
\end{nb}

\begin{nb}[Lie Bracket as a Wriggle]
The Lie bracket describes a motion that cannot be produced by any
single input, but can be produced by a sufficiently fast cyclic
sequence of inputs. For a car-like robot, this is exactly how you
parallel-park: no single input moves the car sideways, but a cyclic
sequence of forward/backward and steering commands does.
\end{nb}

\section{Distributions and Integrability}
\label{sec:distributions}

\begin{definition}[Distribution]
\label{def:distribution}
Let \( C \) be a smooth manifold. A \textbf{distribution} \( \Delta \)
is a map that assigns to each point \( x \in C \) a subspace of the
tangent space \( T_x C \):
\[
\Delta : x \mapsto \Delta(x) \subset T_x C.
\]
It has \textbf{dimension} \( k \), written \( \dim \Delta = k \), if
\( \dim \Delta(x) = k \) for every \( x \in C \).
\end{definition}

For the system \cref{eq:driftless-system} the input distribution is
\( \Delta = \text{span}\{ g_1, \dots, g_m \} \), and it is a
distribution in the sense above whenever the \( g_i \) are
independent at every point.

\begin{definition}[Involutive Distribution]
\label{def:involutive}
A distribution \( \Delta \) is \textbf{involutive} if it is closed
under Lie bracket: for any vector fields \( g, h \) with values in
\( \Delta \), also \( [g, h] \) takes values in \( \Delta \).
\end{definition}

The accessibility distribution \( \Delta_{\mathcal{A}} \) is, by
construction, the smallest involutive distribution containing
\( \Delta \): its \emph{involutive closure}.

\begin{theorem}[Frobenius]
\label{thm:frobenius}
A distribution of constant dimension \( k \) is
\textbf{integrable}, i.e.\ through every point there passes a
\( k \)-dimensional sub-manifold tangent to it, if and only if it is
involutive.
\end{theorem}

Frobenius' theorem is the bridge to \cref{sec:constraints}. Let
\( \Delta = \mathcal{N}(A(q)) \) be the distribution of the admissible
velocities of the Pfaffian constraints \( A(q)\dot q = 0 \). Then:
\begin{itemize}
    \item \( \Delta \) involutive \( \iff \) \( \Delta_{\mathcal{A}} =
          \Delta \) \( \iff \) \( \nu = m \): the constraints are
          completely integrable, hence holonomic, and the motion is
          confined to the integral sub-manifolds of \( \Delta \);
    \item \( \Delta \) not involutive: the constraints are
          nonholonomic. The sub-manifold tangent to
          \( \Delta_{\mathcal{A}} \) (when \( \nu < n \)) is the set
          of accessible configurations.
\end{itemize}
The case \( k = n - 1 \) of \cref{subsec:controllability} is immediate:
a distribution generated by a single vector field is always involutive,
since \( [g, g] = 0 \).

\begin{proposition}
\label{prop:chain-stabilizes}
Let \( \Delta_1 = \Delta \) and
\( \Delta_{i+1} = \Delta_i + \text{span}\{[g,h] : g \in \Delta_1,\
h \in \Delta_i\} \), at points where the dimensions are constant. Then
the chain stabilizes after at most \( n - m \) steps, and
\( \Delta_{\mathcal{A}} = \Delta_{n-m+1} \). In particular, brackets of
length greater than \( n - m + 1 \) are never needed, and
\( \kappa \le n - m + 1 \).
\end{proposition}

\begin{proof}
The chain is increasing, \( \Delta_i \subseteq \Delta_{i+1} \). If
\( \Delta_{i+1} = \Delta_i \), then also \( \Delta_{i+2} = \Delta_{i+1} \),
because the new brackets involve only elements of
\( \Delta_{i+1} = \Delta_i \), which are already in \( \Delta_{i+1} \).
So the chain is either strictly increasing or constant from some point
on. Its dimension starts at \( m \) and cannot exceed \( n \): a strict
increase of at least one dimension per step can happen at most
\( n - m \) times. (Analyticity is not needed for this argument; it is
needed for the converse in \cref{thm:accessibility}.)
\end{proof}

For \( m = 2 \) the bound is \( n - 1 \): the unicycle (\( n = 3 \),
\( \kappa = 2 \)), the bicycle (\( n = 4 \), \( \kappa = 3 \)) and the
chained form (\( \kappa = n - 1 \), see
\cref{prop:chained-controllable}) all attain it.

\section{Exercises}
\label{sec:controllability-exercises}

\begin{enumerate}
    \item \textbf{Heisenberg system.} Let
          \( g_1 = (1, 0, -x_2/2)^\top \) and
          \( g_2 = (0, 1, x_1/2)^\top \) on \( \mathbb{R}^3 \). Compute
          \( [g_1, g_2] \), and decide whether
          \( \dot x = g_1 u_1 + g_2 u_2 \) is controllable. What is
          \( \kappa \)?
    \item \textbf{Unicycle.} Compute \( [g_1, g_2] \) for
          \cref{eq:unicycle} from the Jacobians, and the determinant of
          \( (g_1, g_2, [g_1, g_2]) \).
    \item \textbf{Rear-wheel-drive bicycle.} For \cref{eq:bicycle-rear},
          compute \( g_3 = [g_1, g_2] \), \( g_4 = [g_1, g_3] \), and
          the determinant of \( (g_1, g_2, g_3, g_4) \). For which
          \( \phi \) does the rank drop?
    \item \textbf{Front-wheel-drive bicycle.} Repeat the computation for
          \cref{eq:bicycle-front}. Does the rank drop at
          \( \phi = \pm\pi/2 \)?
    \item \textbf{Chained form.} For the \( (2,4) \) chained form
          compute \( [g_1, g_2] \) and \( [g_1,[g_1,g_2]] \), and verify
          that the rank is 4 and \( \kappa = 3 \).
    \item \textbf{Partial integrability.} Consider on \( \mathbb{R}^4 \)
          the fields \( g_1 = (1, 0, -x_2/2, 0)^\top \),
          \( g_2 = (0, 1, x_1/2, 0)^\top \). Compute \( \nu \) and
          \( m \), classify the constraints with the cases of
          \cref{subsec:controllability}, and say whether the system is
          controllable.
    \item \textbf{Algebra of the bracket.} Prove that
          \( [f, g] = -[g, f] \) and \( [f, f] = 0 \), and that for
          \( f(x) = Ax \), \( g(x) = Bx \) (linear fields) one has
          \( [f, g](x) = (BA - AB)x \).
    \item \textbf{Numerical wriggle.} Simulate the unicycle with the
          sequence of the numerical check in
          \cref{sec:lie-bracket-interpretation}, for
          \( \varepsilon = 0.1, 0.05, 0.025 \). Verify that the final
          lateral displacement scales as \( \varepsilon^2 \).
\end{enumerate}

\paragraph{Solutions (sketch).}
\begin{enumerate}
    \item \( [g_1, g_2] = (0,0,1)^\top \). The matrix
          \( (g_1, g_2, [g_1,g_2]) \) has determinant 1 everywhere, so
          the system is controllable with \( \kappa = 2 \). (This is
          the standard model of a nonholonomic integrator.)
    \item Only \( \partial g_1/\partial q \) is non-zero, with the
          column of \( \theta \) equal to \( (-\sin\theta, \cos\theta,
          0)^\top \). Hence \( [g_1, g_2] = -(\partial g_1/\partial q)
          g_2 = (\sin\theta, -\cos\theta, 0)^\top \), and the
          determinant is \( 1 \) for every \( \theta \).
    \item \( g_3 = -\partial g_1/\partial\phi = (0, 0,
          -1/(l\cos^2\phi), 0)^\top \) and
          \( g_4 = (-\sin\theta, \cos\theta, 0, 0)^\top /
          (l\cos^2\phi) \). The determinant is
          \( 1/(l^2\cos^4\phi) \): the rank drops only at
          \( \phi = \pm\pi/2 \), the singularity of the model.
    \item \( g_3^F = (\cos\theta\sin\phi,\ \sin\theta\sin\phi,\
          -\cos\phi/l,\ 0)^\top \) and
          \( g_4^F = (-\sin\theta,\ \cos\theta,\ 0,\ 0)^\top / l \). In
          the frame aligned with the heading the determinant has
          modulus \( 1/l^2 \), independent of \( \phi \): the rank
          never drops, and \( \kappa = 3 \).
    \item \( g_1 = (1, 0, z_2, z_3)^\top \), \( g_2 = e_2 \),
          \( [g_1, g_2] = -e_3 \), \( [g_1, [g_1, g_2]] = e_4 \). The
          four fields \( g_1, g_2, -e_3, e_4 \) are independent.
    \item \( [g_1, g_2] = (0,0,1,0)^\top = e_3 \), and the next
          brackets vanish: \( \nu = 3 \), \( m = 2 \), \( n = 4 \),
          so \( m < \nu < n \): the constraints are partially
          integrable. The system is not controllable: \( x_4 \) is
          constant (the constraint \( \dot x_4 = 0 \) is integrable).
    \item Antisymmetry is immediate from \( J_g f - J_f g \). For the
          linear fields, \( J_f = A \), \( J_g = B \), so
          \( [f,g] = BAx - ABx \).
    \item The lateral displacement is \( -\varepsilon\sin\varepsilon
          \approx -\varepsilon^2 \): halving \( \varepsilon \) divides
          it by about 4.
\end{enumerate}

\section{Takeaways}
\label{sec:controllability-takeaways}

\begin{enumerate}
    \item For driftless systems the accessibility rank condition is
          necessary and sufficient for controllability
          (Chow--Rashevskii); with drift it only guarantees
          accessibility.
    \item The Lie bracket is bilinear, antisymmetric, and satisfies the
          Jacobi identity. A cyclic input sequence of duration
          \( 4\varepsilon \) produces a displacement
          \( \varepsilon^2 [g_1, g_2] + O(\varepsilon^3) \) (a
          \textbf{wriggle}).
    \item Frobenius: a constant-dimension distribution is integrable if
          and only if it is involutive. Holonomic constraints
          correspond to \( \Delta_{\mathcal{A}} = \Delta \).
    \item The chain of brackets stabilizes after at most \( n - m \)
          steps, so \( \kappa \le n - m + 1 \): the unicycle, the
          bicycle and the chained form (\cref{ch:wmr-mechanics}) all
          attain the bound for \( m = 2 \).
\end{enumerate}